Выпуклые задачи --- класс, в котором локальный минимум всегда глобальный. Сначала вводятся выпуклые множества и операции, сохраняющие выпуклость, затем выпуклые функции и два их критерия: через касательную (\hyperref[note:18BD8497-CEBD-4962-8144-474B489606C3]{дифференциальный критерий}) и через гессиан (\hyperref[note:86FCE16B-C82E-4416-851C-0E0AD0FDA382]{критерий вторых производных}). Главное следствие для оптимизации --- \hyperref[note:035B8935-65A9-4BB4-BDA3-60A60FF278F0]{минимум выпуклой функции достигается в точке с нулевым градиентом}.
